% !TEX root = ../main.tex
\chapter{Discussion and Analysis}
\label{ch:evaluation}

This chapter discusses the experimental results, compares them to existing literature, and outlines the limitations of the study. The findings demonstrate that target transformations are key to stabilizing training in imagination-based reinforcement learning.

%--------------------------------------------------------------
\section{Interpretation of Results}
%--------------------------------------------------------------

\subsection{Effect of Imagination-Based Training}

The comparison between Condition~A (model-free actor-critic) and Condition~B (imagination-based actor-critic with MSE losses) illustrates the trade-offs of learning from simulated experience. In CartPole-v1, Condition~B showed higher sample efficiency in early training than Condition~A, reaching near-optimal returns in fewer environment interactions (see \cref{fig:cartpole-learning}). This outcome is consistent with model-based reinforcement learning theory: generating simulated rollouts from a trained world model allows the agent to reuse transitions, reducing the sample complexity of learning.

However, model-based training introduces prediction error from the world model into the optimization process. When training on simulated trajectories, errors in the world model are propagated to the actor and critic updates. In CartPole-v1, which features a constant $+1$ reward and smooth state transitions, these compounding errors remain small. As shown in \cref{fig:cartpole-loss}, the critic loss for Condition~B remained bounded.

In contrast, in LunarLander-v3, Condition~B failed to learn and experienced policy collapse, with mean returns dropping to approximately $-1000$ (\cref{fig:lunar-learning}). This indicates that imagination-based training without proper target scaling can destabilize policy learning.

\subsection{The Value Explosion Problem}

The failure of Condition~B in LunarLander-v3 is associated with value function divergence. As shown in \cref{fig:lunar-loss}, the critic loss for Condition~B escalated to approximately $10^{15}$, which corrupted the value estimates and the policy gradient signal.

This divergence is driven by the compounding of prediction errors across imagined rollouts. During imagination, the agent generates trajectories of horizon $H$ by autoregressively applying the world model. At step $t$, the predicted next state contains prediction errors. Because subsequent steps are conditioned on previous predictions, these errors accumulate. Over an $H$-step horizon, the cumulative prediction error grows.

In LunarLander-v3, this accumulation is amplified by the reward structure, which includes terminal rewards of $-100$ and $+100$ along with continuous shaping rewards. When the critic fits these variable target values using standard MSE loss, prediction errors result in large squared residuals. For example, a difference of $200$ between the target and predicted values yields a squared error of $40{,}000$ before temporal compounding is factored in.

The resulting large losses produce large gradients, which push the critic's parameters to produce extreme predictions. This feedback loop (where prediction errors yield large gradients that increase subsequent prediction errors) leads to divergence. This instability motivated the development of symlog losses in DreamerV3~\citep{hafner2023dreamerv3}.

\subsection{Symlog as a Stabilizer}

Condition~C demonstrates that symlog-transformed losses prevent value divergence. In LunarLander-v3, the critic loss under Condition~C remained stable throughout training (\cref{fig:lunar-loss}), while the mean returns stabilized at approximately $-140$, matching the model-free baseline (\cref{fig:lunar-learning}).

The symlog function, defined as:
\begin{equation}
    \operatorname{symlog}(x) = \operatorname{sign}(x) \ln(|x| + 1)
\end{equation}
compresses values logarithmically while preserving their sign. A target of $-100$ maps to $\operatorname{symlog}(-100) \approx -4.62$, and $+100$ maps to $\approx +4.62$. The dynamic range is thus compressed.

This compression bounds the squared residuals. The maximum squared error in symlog space is restricted to approximately $(4.62 - (-4.62))^2 \approx 85.4$, compared to values up to $40{,}000$ in the original space. Bounding the loss values constrains the gradient magnitudes, preventing the feedback loop that drives divergence.

Additionally, the world model's reward predictor benefits from the symlog transformation. As shown in \cref{fig:lunar-wm}, the reward prediction error under Condition~C is lower than that of Condition~B. Predicting in the compressed symlog space simplifies regression because the target values are bounded.

%--------------------------------------------------------------
\section{Comparison with Existing Literature}
%--------------------------------------------------------------

These results support the findings of \citet{hafner2023dreamerv3} regarding the role of symlog transformations in stabilizing world-model-based learning. While DreamerV3 combines symlog predictions with several other components (two-hot encoding, KL balancing, and the RSSM), our study isolates the effect of symlog. By using raw state-space observations, a feedforward MLP world model, and standard actor-critic updates, we verify that the symlog transformation alone prevents value divergence.

Our hybrid Dyna-style training approach, which uses both real and simulated transitions, connects to classic planning architectures. Interleaving model-free updates on real data with planning on simulated data stabilizes learning when the world model is still inaccurate early in training. Real transitions provide a baseline that helps preserve policy quality.

%--------------------------------------------------------------
\section{Limitations}
%--------------------------------------------------------------

Several limitations of this study should be noted:

\begin{itemize}
    \item \textbf{Simplified world model architecture:} The feedforward MLP world model operating in the raw observation space is simpler than recurrent latent-space models like the RSSM in DreamerV3. This simplification isolates the effect of symlog but limits the capacity to capture long-term dependencies.
    \item \textbf{Discrete action spaces:} We evaluated the configurations on environments with discrete action spaces. The interaction between symlog transformations and continuous action spaces remains to be evaluated.
    \item \textbf{Limited environment diversity:} The study is restricted to CartPole-v1 and LunarLander-v3. Evaluating a wider range of environments with diverse reward scales would strengthen the findings.
    \item \textbf{Omission of two-hot encoding:} We did not evaluate two-hot categorical encoding, which DreamerV3 combines with symlog.
    \item \textbf{Incomplete task performance:} While Condition~C prevented value explosion in LunarLander-v3, the final returns remained around $-140$, falling short of the solved threshold of $+200$. This indicates that while symlog stabilizes training, additional architectural or algorithmic components are needed for high performance.
\end{itemize}

%--------------------------------------------------------------
\section{Summary}
%--------------------------------------------------------------

This chapter analyzed the experimental results, focusing on error compounding and the stabilizing properties of the symlog transformation. The symlog function prevents value divergence by bounding the dynamic range of the regression targets. Our results isolate the contribution of symlog in a simplified setting. The limitations identified here suggest directions for future work, which are discussed in Chapter~\ref{ch:con}.